% ============================================================================
\chapter{费马大定理}
\label{ch:flt}
% ============================================================================

% ----------------------------------------------------------------------------
\section{费马大定理的陈述}
% ----------------------------------------------------------------------------

\begin{theorem}[费马大定理, Wiles 1995]
\label{thm:flt}
方程
\[
  x^n + y^n = z^n
\]
在 $n \geq 3$ 时没有正整数解。
\end{theorem}

% ----------------------------------------------------------------------------
\section{历史背景}
% ----------------------------------------------------------------------------

% TODO:
% - Fermat 的批注（1637）
% - 特殊情形：n=3 (Euler), n=4 (Fermat), n=5 (Dirichlet-Legendre), 
%   n=7 (Lamé), 正则素数 (Kummer)
% - 计算机验证的进展

% ----------------------------------------------------------------------------
\section{Frey 曲线}
% ----------------------------------------------------------------------------

% TODO:
% - 若 aᵖ + bᵖ = cᵖ，构造 Frey 曲线 E: y² = x(x-aᵖ)(x+bᵖ)
% - Frey 曲线的性质：半稳定，导子 N = rad(abc)²

% ----------------------------------------------------------------------------
\section{Ribet 定理（Level Lowering）}
% ----------------------------------------------------------------------------

\begin{theorem}[Ribet, 1990]
设 $E/\Q$ 为半稳定椭圆曲线。若 $E$ 模性（即存在 $f \in \Sk_2(\GammaO(N_E))$
使得 $\rho_{E,p} \cong \rho_{f,p}$），则可以降低 level：
存在 $g \in \Sk_2(\GammaO(2))$ 使得 $\rho_{E,p} \cong \rho_{g,p}$。
\end{theorem}

% TODO: 
% - 但 S₂(Γ₀(2)) = 0，矛盾！
% - 因此 Frey 曲线不存在 → FLT 成立

% ----------------------------------------------------------------------------
\section{完整的证明逻辑链}
% ----------------------------------------------------------------------------

% TODO:
% 假设 aᵖ + bᵖ = cᵖ 有正整数解（p ≥ 5 为奇素数）
% → 构造 Frey 曲线 E
% → E 半稳定（Frey）
% → E 模性（Wiles-Taylor 的模性定理）
% → Level lowering 到 Γ₀(2)（Ribet 定理）
% → S₂(Γ₀(2)) = 0，矛盾
% → FLT 成立

% ----------------------------------------------------------------------------
\section{后续发展}
% ----------------------------------------------------------------------------

% TODO:
% - 完全模性定理（Breuil-Conrad-Diamond-Taylor, 2001）
% - Serre 猜想（Khare-Wintenberger, 2009）
% - 推广到全实域（Freitas-Le Hung-Siksek, 2015）

% ----------------------------------------------------------------------------
\section{习题}
% ----------------------------------------------------------------------------
